
\section{Distâncias e Propriedades Geométricas no Espaço de Representação}
\label{sec:representation_geometry}

O pipeline núcleo-transformação associa a um conjunto finito de objetos uma
matriz simétrica semidefinida positiva
\begin{equation}\label{eq:transformed_gram_geometry}
    \mathbf{G}
    =
    T(\widetilde{\mathbf{G}})
    \in
    \mathbb{S}_+^n.
\end{equation}
Essa matriz determina uma geometria empírica sobre as representações
analisadas. De fato, como $\mathbf{G}$ é semidefinida positiva, existem um
espaço de Hilbert $\mathcal{E}$ e vetores
$\varphi_1,\dots,\varphi_n\in\mathcal{E}$ tais que
\begin{equation}\label{eq:empirical_feature_realization}
    G_{ij}
    =
    \langle
        \varphi_i,
        \varphi_j
    \rangle_{\mathcal{E}}.
\end{equation}

A representação explícita dos vetores $\varphi_i$ não é necessária. Seus
produtos internos, normas e distâncias podem ser obtidos diretamente a
partir de $\mathbf{G}$. Desse modo, objetos pertencentes a espaços de
representação possivelmente heterogêneos ou sem uma métrica previamente
especificada passam a admitir uma geometria mensurável no espaço de
características induzido pelo pipeline.

Essa geometria não deve ser interpretada como uma propriedade intrínseca e
única do domínio original. Ela depende da função núcleo e da transformação
pós-\textit{kernel} adotadas. Diferentes pipelines podem, portanto, induzir
diferentes noções de proximidade entre os mesmos objetos, de acordo com as
propriedades geométricas preservadas ou descartadas em cada construção.

\subsection{Pseudométrica induzida pela matriz de Gram}
\label{subsec:gram_pseudometric}

A distância entre as representações dos objetos $i$ e $j$ é definida pela
norma da diferença no espaço de características:
\begin{equation}\label{eq:induced_distance}
    d_{\mathbf{G}}(i,j)
    =
    \|
        \varphi_i-\varphi_j
    \|_{\mathcal{E}}.
\end{equation}
Expandindo a norma e utilizando \eqref{eq:empirical_feature_realization},
obtém-se
\begin{equation}\label{eq:gram_distance}
    d_{\mathbf{G}}^2(i,j)
    =
    G_{ii}
    +
    G_{jj}
    -
    2G_{ij}.
\end{equation}

A matriz de distâncias associada a $\mathbf{G}$ é
\begin{equation}\label{eq:distance_matrix}
    \mathbf{D}_{\mathbf{G}}
    =
    \left(
        d_{\mathbf{G}}(i,j)
    \right)_{i,j=1}^{n}.
\end{equation}
Por construção, $\mathbf{D}_{\mathbf{G}}$ é simétrica, possui diagonal
nula e satisfaz a desigualdade triangular.

Em geral, $d_{\mathbf{G}}$ é uma pseudométrica sobre o conjunto de índices
dos objetos. É possível que dois objetos distintos satisfaçam
\begin{equation}\label{eq:zero_pseudodistance}
    d_{\mathbf{G}}(i,j)
    =
    0,
    \qquad
    i\neq j,
\end{equation}
quando suas representações coincidem no espaço induzido:
\begin{equation*}
    \varphi_i
    =
    \varphi_j.
\end{equation*}
Essa possibilidade não constitui uma deficiência da construção. Ela
expressa que o pipeline considera os dois objetos indistinguíveis segundo
a noção de similaridade adotada.

\begin{proposition}[Propriedades da distância induzida]
\label{prop:induced_pseudometric}
A função
\begin{equation*}
    d_{\mathbf{G}}:
    \{1,\dots,n\}
    \times
    \{1,\dots,n\}
    \longrightarrow
    \mathbb{R}_{\geq 0}
\end{equation*}
definida em \eqref{eq:induced_distance} satisfaz:
\begin{enumerate}
    \item $d_{\mathbf{G}}(i,j)\geq 0$,
    \item $d_{\mathbf{G}}(i,j)=d_{\mathbf{G}}(j,i)$,
    \item $d_{\mathbf{G}}(i,i)=0$,
    \item $d_{\mathbf{G}}(i,k) \leq d_{\mathbf{G}}(i,j) + d_{\mathbf{G}}(j,k)$.

\end{enumerate}
Além disso,
\begin{equation}
    d_{\mathbf{G}}(i,j)
    =
    0
    \iff
    \varphi_i
    =
    \varphi_j.
\end{equation}
Consequentemente, $d_{\mathbf{G}}$ é uma pseudométrica sobre os objetos e
uma métrica sobre as classes de equivalência induzidas por representações
coincidentes.
\end{proposition}

\begin{proof}
As propriedades de não negatividade, simetria e desigualdade triangular
decorrem diretamente do fato de que $d_{\mathbf{G}}$ é obtida pela norma
de um espaço de Hilbert. Além disso,
\begin{equation*}
    d_{\mathbf{G}}(i,j)
    =
    0
\end{equation*}
se e somente se
\begin{equation*}
    \|
        \varphi_i-\varphi_j
    \|_{\mathcal{E}}
    =
    0,
\end{equation*}
o que ocorre se e somente se
$\varphi_i=\varphi_j$.
\end{proof}

\subsection{Equivalência representacional}
\label{subsec:representational_equivalence}

A pseudométrica induz naturalmente a relação
\begin{equation}\label{eq:representation_equivalence}
    i\sim_{\mathbf{G}}j
    \iff
    d_{\mathbf{G}}(i,j)=0.
\end{equation}
Equivalentemente,
\begin{equation*}
    i\sim_{\mathbf{G}}j
    \iff
    \varphi_i=\varphi_j.
\end{equation*}
Essa relação é reflexiva, simétrica e transitiva e, portanto, particiona o
conjunto de objetos em classes de equivalência representacional.

Objetos pertencentes à mesma classe não são necessariamente iguais no
domínio original. Eles são apenas indistinguíveis segundo o espaço de
representação, o núcleo e a transformação pós-\textit{kernel} escolhidos.
A equivalência é, portanto, relativa ao pipeline:
\begin{equation*}
    \sim_{\mathbf{G}}
    =
    \sim_{(\kappa,T)}.
\end{equation*}

O espaço quociente
\begin{equation}
    \{1,\dots,n\}/{\sim_{\mathbf{G}}}
\end{equation}
reúne as classes de objetos com representação coincidente. Sobre esse
quociente, $d_{\mathbf{G}}$ é uma métrica propriamente dita.

\subsection{Geometria relativa da amostra}
\label{subsec:sample_relative_geometry}

A matriz $\mathbf{D}_{\mathbf{G}}$ reúne a geometria métrica relativa dos
objetos observados. Suas entradas descrevem como as representações se
organizam umas em relação às outras, sem exigir acesso às suas coordenadas
explícitas no espaço de características. Essa estrutura permite analisar
proximidade, equivalência representacional, vizinhanças, isolamento local e
agrupamentos, sem exigir que a pseudométrica separe todos os objetos.

\subsubsection*{Vizinhanças induzidas}

Para um objeto de referência $i_0$, a vizinhança de raio
$\varepsilon\geq 0$ é definida por
\begin{equation}\label{eq:epsilon_neighborhood}
    \mathcal{N}_{\varepsilon}(i_0)
    =
    \left\{
        j\in\{1,\dots,n\}
        :
        d_{\mathbf{G}}(i_0,j)
        \leq
        \varepsilon
    \right\}.
\end{equation}
Alternativamente, pode-se considerar a vizinhança
$\mathcal{N}_k(i_0)$ formada pelos $k$ objetos mais próximos de $i_0$,
com respeito a $d_{\mathbf{G}}$.

Essas vizinhanças descrevem a organização local da amostra. Objetos com
muitos vizinhos a pequenas distâncias ocupam regiões densas no espaço
induzido, enquanto objetos cujos vizinhos mais próximos permanecem
distantes ocupam regiões relativamente isoladas. Como
$d_{\mathbf{G}}$ é uma pseudométrica, uma vizinhança de raio zero pode
conter mais de um objeto:
\begin{equation}
    \mathcal{N}_{0}(i_0)
    =
    \left\{
        j:
        j\sim_{\mathbf{G}} i_0
    \right\}.
\end{equation}
Assim, $\mathcal{N}_{0}(i_0)$ coincide com a classe de equivalência
representacional de $i_0$.

\subsubsection*{Conjuntos de nível da distância}

Fixado um objeto de referência $i_0$, define-se a função distância
\begin{equation}\label{eq:sample_distance_function}
    p_{i_0}(j)
    =
    d_{\mathbf{G}}(i_0,j).
\end{equation}
Para cada valor $r\geq 0$, o conjunto de nível associado a $i_0$ é
\begin{equation}\label{eq:sample_level_set}
    \Gamma_r(i_0)
    =
    \left\{
        j\in\{1,\dots,n\}
        :
        p_{i_0}(j)=r
    \right\}.
\end{equation}
Os elementos de $\Gamma_r(i_0)$ possuem representações equidistantes da
representação de referência. À medida que $r$ varia, esses conjuntos
organizam a amostra em camadas de equidistância ao redor de $i_0$.

No caso $r=0$, tem-se
\begin{equation}
    \Gamma_0(i_0)
    =
    \left\{
        j:
        d_{\mathbf{G}}(i_0,j)=0
    \right\}
    =
    [i_0]_{\sim_{\mathbf{G}}},
\end{equation}
de modo que o conjunto de nível nulo coincide com a classe de equivalência
representacional de $i_0$. Para $r>0$, os conjuntos de nível descrevem as
diferentes escalas de afastamento em relação à referência.

Como o conjunto de objetos é finito, $\Gamma_r(i_0)$ pode ser vazio para
valores de $r$ que não ocorram entre as entradas da matriz de distâncias.
Nesse regime, os conjuntos de nível devem ser entendidos como camadas
discretas da amostra, e não necessariamente como curvas geométricas
contínuas.

\subsubsection*{Faixas de nível}

Em aplicações empíricas, distâncias exatamente iguais podem ser raras. Para
obter uma descrição mais estável, podem ser consideradas faixas de nível
com espessura $\delta>0$:
\begin{equation}\label{eq:sample_level_band}
    \Gamma_{r,\delta}(i_0)
    =
    \left\{
        j\in\{1,\dots,n\}
        :
        \left|
            d_{\mathbf{G}}(i_0,j)-r
        \right|
        \leq\delta
    \right\}.
\end{equation}
Essas faixas agrupam objetos cujas distâncias à referência pertencem a uma
mesma escala e constituem uma aproximação empírica dos conjuntos de nível.

Quando os objetos estão associados a posições de um domínio espacial ou
contínuo, a função distância pode ser avaliada ou interpolada sobre esse
domínio. Nesse caso, os conjuntos de nível assumem a forma de curvas ou
superfícies de nível, permitindo visualizar a organização geométrica da
representação em torno de uma referência. A interpretação dessas curvas
dependerá da representação e do pipeline empregados na aplicação.

O Algoritmo \ref{alg:geometria} resume o cálculo dessas quantidades a partir
da matriz de Gram, sem que as representações $\varphi_i$ precisem ser obtidas.

\begin{algorithm}[htbp]
\caption{Relative geometry induced by the Gram matrix}
\label{alg:geometria}
\begin{algorithmic}[1]
\Require $\mathbf{G} \in \mathbb{S}_+^n$, reference object $i_0$, number of neighbors $k$, radius $r$ and thickness $\delta$
\Ensure $\mathbf{D}_{\mathbf{G}}$, $\mathcal{N}_k(i_0)$ and $\Gamma_{r,\delta}(i_0)$
\For{$i, j = 1,\dots,n$}
    \State $(\mathbf{D}_{\mathbf{G}})_{ij} \gets \sqrt{\max\left(G_{ii} + G_{jj} - 2G_{ij},\, 0\right)}$ \Comment{Equation \eqref{eq:gram_distance}}
\EndFor
\State $\mathcal{N}_k(i_0) \gets$ the $k$ objects closest to $i_0$ according to $\mathbf{D}_{\mathbf{G}}$
\State $\Gamma_{r,\delta}(i_0) \gets \{\, j : |(\mathbf{D}_{\mathbf{G}})_{i_0 j} - r| \leq \delta \,\}$
\end{algorithmic}
\end{algorithm}
\subsection{Custo computacional}\label{subsec:custo}

Esta subseção analisa o custo assintótico dos Algoritmos \ref{alg:pipeline}, \ref{alg:geometria} e \ref{alg:energia}, isto é, do método em si, independentemente da instanciação. O custo de obter as fontes, que depende do modelo analisado, é discutido na Subseção \ref{subsec:reprodutibilidade}. Utiliza-se a notação usual \citep{Cormen2009}: $O(g)$ indica um limite superior, $\Omega(g)$ um limite inferior e $\Theta(g)$ um limite justo, que vale como superior e inferior ao mesmo tempo. Sejam $n$ o número de objetos, $d$ o número de fontes, $c_l$ o custo de avaliar o \textit{kernel} $\kappa_l$ sobre um par de objetos, $r$ o número de modos retidos e $k$ o número de vizinhos. Com \textit{kernels} lineares, angulares ou gaussianos sobre representações de dimensão $D_l$, tem-se $c_l = \Theta(D_l)$, e denota-se $D = \sum_l D_l$.

Veja que os três algoritmos são determinísticos e não tomam decisões que dependam dos valores dos dados: para um mesmo $n$ e uma mesma dimensão das fontes, eles executam sempre as mesmas operações. Por esse motivo, o melhor caso, o pior caso e o caso médio coincidem, e o custo de cada algoritmo pode ser descrito por um limite justo $\Theta$. O limite inferior $\Omega$ é utilizado com outro sentido, o do \textit{problema}: ele indica o mínimo que qualquer algoritmo precisa gastar para produzir a mesma saída, e a distância entre ele e o custo do algoritmo mostra onde há espaço para melhorias.

\textbf{Construção da matriz de Gram.} O Algoritmo \ref{alg:pipeline} avalia o \textit{kernel} de cada fonte sobre os $n(n+1)/2$ pares distintos, o que custa $\Theta(n^2 c_l)$ por fonte e $\Theta(n^2 D)$ no total. As transformações por componente utilizadas neste trabalho, a normalização angular $T_a$, o deslocamento por $\mathbf{J}$, o centramento $T_c$ e a normalização em traço $T_{\mathrm{tr}}$, operam entrada a entrada ou por somas de linhas e colunas, e custam $\Theta(n^2)$ cada uma. Veja que o centramento não exige os produtos $H\widetilde{\mathbf{G}}H$, visto que $(H\widetilde{\mathbf{G}}H)_{ij} = \widetilde{G}_{ij} - \bar{g}_i - \bar{g}_j + \bar{g}$, em que $\bar{g}_i$ é a média da linha $i$ e $\bar{g}$ a média de todas as entradas. O produto de Hadamard das $d$ componentes custa $\Theta(d\,n^2)$. Sendo assim, com transformações entrada a entrada, a construção custa $\Theta(n^2 D)$ em tempo e $\Theta(n^2)$ em memória, visto que as fontes podem ser processadas uma de cada vez. As transformações espectrais da Equação \eqref{eq:spectral_transformation}, como $T_w$, exigem uma decomposição espectral e custam $\Theta(n^3)$ \citep{GolubVanLoan2013}. Qualquer algoritmo que produza $\mathbf{G}$ precisa escrever as suas $n(n+1)/2$ entradas e ler as $nD$ coordenadas das fontes, de modo que o problema tem limite inferior $\Omega(n^2 + nD)$. A distância entre esse limite e $\Theta(n^2 D)$ é exatamente o custo do produto $XX^\top$ das representações, que algoritmos rápidos de multiplicação de matrizes reduzem apenas assintoticamente.

\textbf{Geometria.} No Algoritmo \ref{alg:geometria}, cada entrada da matriz de distâncias é obtida da Equação \eqref{eq:gram_distance} em tempo constante, e a matriz inteira custa $\Theta(n^2)$. Como a saída tem $n^2$ entradas, esse custo coincide com o limite inferior $\Omega(n^2)$, isto é, a construção da geometria é ótima. A vizinhança $\mathcal{N}_k(i_0)$ e a faixa de nível $\Gamma_{r,\delta}(i_0)$ de uma referência exigem apenas a linha $i_0$ da matriz e custam $\Theta(n)$, a primeira por seleção em tempo linear \citep{Cormen2009}, e ordenar a linha, como no algoritmo, custa $\Theta(n \log n)$. As leituras que tomam cada objeto como referência, como os $k$-vizinhos \textit{leave-one-out}, custam $\Theta(n^2)$ por seleção e $\Theta(n^2 \log n)$ por ordenação.

\textbf{Energia espectral.} O Algoritmo \ref{alg:energia} realiza a decomposição espectral completa de $\mathbf{G}$, que custa $\Theta(n^3)$ com os algoritmos densos usuais para matrizes simétricas \citep{GolubVanLoan2013}, e a partir dela obtém as energias e a razão $\widehat{E}_{\mathrm{ratio}}$ de todos os objetos em $\Theta(nr)$. Essa é a etapa mais cara do método, mas ela pode ser evitada em duas partes. Primeiro, a razão de participação da Equação \eqref{eq:r_participation} não exige os autovalores, visto que, para uma matriz simétrica, $\sum_k \hat\lambda_k = \operatorname{tr}(\mathbf{G})$ e $\sum_k \hat\lambda_k^2 = \|\mathbf{G}\|_F^2$, de modo que $r$ é obtido em $\Theta(n^2)$. Segundo, as leituras utilizam apenas os $r$ modos retidos, e estes podem ser obtidos por métodos de Lanczos ou aleatorizados em $O(n^2 r)$ \citep{Halko2011}. O problema tem limite inferior $\Omega(n^2)$, visto que $\mathbf{G}$ precisa ser lida.

A Tabela \ref{tab:custo} resume a análise. Com a decomposição espectral completa, o método custa $\Theta(n^2 D + n^3)$ por matriz de Gram, e com a decomposição truncada custa $O\bigl(n^2 (D + r)\bigr)$, contra um limite inferior de $\Omega(n^2 + nD)$. Veja que, em ambos os casos, o custo é polinomial em $n$ e linear na dimensão das fontes, e não depende do número de atributos da entrada nem do número de parâmetros do modelo, que entram apenas na obtenção das fontes.

\begin{table}[htbp]
\centering
\caption{Custo assintótico das etapas do método, para $n$ objetos, $d$ fontes de dimensão total $D$, $r$ modos retidos e $k$ vizinhos. A segunda coluna é o custo dos algoritmos como descritos, a terceira é uma alternativa mais barata, quando existe, e a quarta é o limite inferior do problema. As células vazias indicam que não há alternativa mais barata.}
\label{tab:custo}
\small
\begin{tabular}{l c c c}
\hline
Etapa & Algoritmo & Alternativa & Problema \\
\hline
Gram bruta de todas as fontes & $\Theta(n^2 D)$ & & $\Omega(n^2 + nD)$ \\
$T_a$, deslocamento, $T_c$, $T_{\mathrm{tr}}$, Hadamard & $\Theta(d\,n^2)$ & & $\Omega(n^2)$ \\
Matriz de distâncias & $\Theta(n^2)$ & & $\Omega(n^2)$ \\
$\mathcal{N}_k(i_0)$ e $\Gamma_{r,\delta}(i_0)$ & $\Theta(n \log n)$ & $\Theta(n)$ & $\Omega(n)$ \\
$k$-vizinhos de todos os objetos & $\Theta(n^2 \log n)$ & $\Theta(n^2)$ & $\Omega(n^2)$ \\
Razão de participação & $\Theta(n^3)$ & $\Theta(n^2)$ & $\Omega(n^2)$ \\
Decomposição espectral & $\Theta(n^3)$ & $O(n^2 r)$ & $\Omega(n^2)$ \\
Energias e $\widehat{E}_{\mathrm{ratio}}$, dada a decomposição & $\Theta(nr)$ & & $\Omega(nr)$ \\
\hline
Total & $\Theta(n^2 D + n^3)$ & $O\bigl(n^2(D + r)\bigr)$ & $\Omega(n^2 + nD)$ \\
\hline
\end{tabular}
\end{table}
